\histitem{Representations of locally compact groups}

Until 1939, the study of representations of topological groups was limited to the compact case and the commutative case. Physicists, for their part, had found in quantum mechanics many reasons to study irreducible unitary representations in order to seek Hilbert spaces capable of modeling the states of elementary particles. P. Dirac, building on work by E. Majorana, had pointed out around 1936 the interest, in this perspective, of the Lorentz group SO(3,1) and the Poincaré group SO(3,1) × \(R^{4}\). Their methods remained far removed from those of the mathematicians of the time. But in 1939, E. Wigner {[}90{]} opened a breach when he classified the irreducible unitary representations of the Poincaré group, building on the work of Murray and von Neumann.

It is to Gelfand that we owe, it seems, having glimpsed a general path toward the study of representations of locally compact groups. In a memoir written in 1942 (cf. {[}23{]}, vol. II, p. 3--17), he notes with Raikov that such a theory is made possible by the link between unitary representations and functions of positive type. They prove the existence of a Hilbertian realization for every function of positive type (V, p. 432, th. 1) and deduce from it that irreducible unitary representations separate the points of G (V, p. 454, th. 4). The positive linear forms on the involutive algebra L \(^{1}\) (G) play an essential role: Gelfand and Raikov show the link that connects them to functions of positive type (cf. V, p. 448, prop. 13). The enveloping C*-algebra of an involutive algebra is presented (without a name) in {[}24, § 48{]}, although it is not yet directly a question of associating with G the algebra Stell(G) (cf. 9, def. of I, p. 125).

After 1945, the study of unitary representations developed by leaps and bounds. The interaction with C*-algebras, as well as with the analysis of the foundations of quantum theory, quickly provided valuable general results: around 1947, Segal associated with a locally compact group various \(^{(9)}\) stellar algebras {[}67{]} and wove the link between their representations and those of the group under consideration. Shortly thereafter, G. Mackey extracts the notion of induced representation {[}42{]} which will be omnipresent in the concrete results on unitary representations.

At the same time, the theory is considerably stimulated by the study of examples that reveal its variety and depth. Gelfand and Naimark, studying the example of SL(2, C), isolated in 1947 a special case of the notion of induced representation, then showed that it is the key to the study of the representations of this group {[}23, vol. II, p. 41--124{]}. In the same year, V. Bargmann {[}5{]} studied the group SL(2, R) and discovered the existence of a countable family of square-integrable representations, called the "discrete series," as well as the orthogonality relations between their matrix coefficients (cf. prop. 8 of V, p. 424). As Godement immediately realized {[}25{]}, the Bargmann orthogonality relations are satisfied by all square-integrable representations of locally compact groups.

Soon the study of particular classes of groups would again, and for a long time, take precedence over abstract theory. Harish-Chandra, in particular, would undertake the general study of representations of reductive Lie groups --- an epic task that we cannot describe here, any more than the prodigious consequences that Langlands would foresee for number theory.

